% Options for packages loaded elsewhere
\PassOptionsToPackage{unicode}{hyperref}
\PassOptionsToPackage{hyphens}{url}
\PassOptionsToPackage{dvipsnames,svgnames,x11names}{xcolor}
%
\documentclass[
  11pt,
  a4paper,
]{ltjsarticle}
\usepackage{amsmath,amssymb}
\usepackage{iftex}
\ifPDFTeX
  \usepackage[T1]{fontenc}
  \usepackage[utf8]{inputenc}
  \usepackage{textcomp} % provide euro and other symbols
\else % if luatex or xetex
  \usepackage{unicode-math} % this also loads fontspec
  \defaultfontfeatures{Scale=MatchLowercase}
  \defaultfontfeatures[\rmfamily]{Ligatures=TeX,Scale=1}
\fi
\usepackage{lmodern}
\ifPDFTeX\else
  % xetex/luatex font selection
\fi
% Use upquote if available, for straight quotes in verbatim environments
\IfFileExists{upquote.sty}{\usepackage{upquote}}{}
\IfFileExists{microtype.sty}{% use microtype if available
  \usepackage[]{microtype}
  \UseMicrotypeSet[protrusion]{basicmath} % disable protrusion for tt fonts
}{}
\makeatletter
\@ifundefined{KOMAClassName}{% if non-KOMA class
  \IfFileExists{parskip.sty}{%
    \usepackage{parskip}
  }{% else
    \setlength{\parindent}{0pt}
    \setlength{\parskip}{6pt plus 2pt minus 1pt}}
}{% if KOMA class
  \KOMAoptions{parskip=half}}
\makeatother
\usepackage{xcolor}
\usepackage[margin=24mm]{geometry}
\setlength{\emergencystretch}{3em} % prevent overfull lines
\providecommand{\tightlist}{%
  \setlength{\itemsep}{0pt}\setlength{\parskip}{0pt}}
\setcounter{secnumdepth}{5}
\ifLuaTeX
  \usepackage{selnolig}  % disable illegal ligatures
\fi
\IfFileExists{bookmark.sty}{\usepackage{bookmark}}{\usepackage{hyperref}}
\IfFileExists{xurl.sty}{\usepackage{xurl}}{} % add URL line breaks if available
\urlstyle{same}
\hypersetup{
  colorlinks=true,
  linkcolor={blue},
  filecolor={Maroon},
  citecolor={Blue},
  urlcolor={blue},
  pdfcreator={LaTeX via pandoc}}

\author{}
\date{}

\begin{document}

\hypertarget{ux6f14ux7fd2-03-ux9023ux7acbux4e00ux6b21ux65b9ux7a0bux5f0f}{%
\section{演習 03 ---
連立一次方程式}\label{ux6f14ux7fd2-03-ux9023ux7acbux4e00ux6b21ux65b9ux7a0bux5f0f}}

\hypertarget{gauss-ux6d88ux53bb}{%
\subsection{1. Gauss 消去}\label{gauss-ux6d88ux53bb}}

\[
\begin{cases}
x+2y-z=1\\
2x+4y+z=4\\
-x-2y+2z=1
\end{cases}
\]

を Gauss 消去で解いてください。

\hypertarget{ux89e3ux306eux5b58ux5728}{%
\subsection{2. 解の存在}\label{ux89e3ux306eux5b58ux5728}}

\texttt{Ax=b} が解を持つことと \texttt{b∈Col(A)}
が同値であることを、行列積を列の一次結合として説明してください。

\hypertarget{ux89e3ux96c6ux5408}{%
\subsection{3. 解集合}\label{ux89e3ux96c6ux5408}}

\texttt{x\_p} が \texttt{Ax=b} の一つの解であるとき、全解が
\texttt{x\_p+ker(A)} と書けることを証明してください。

\hypertarget{ux4e00ux610fux6027}{%
\subsection{4. 一意性}\label{ux4e00ux610fux6027}}

\texttt{Ax=b} の解が存在すると仮定します。解が一意であることと
\texttt{ker(A)=\{0\}} が同値であることを示してください。

\hypertarget{ux9006ux884cux5217}{%
\subsection{5. 逆行列}\label{ux9006ux884cux5217}}

「逆行列を持つ」「full rank」「零空間が
\texttt{\{0\}}」の関係を説明してください。

\end{document}
